%%=====================================================================
\section{The quaternionic fibre and the corner-transit parity}\label{sec:parity}
%%=====================================================================

This section supplements the arithmetic of
Section~\ref{sec:arithmetic} with a group-theoretic structure carried
by the quaternion group, and establishes three results: an orbit
decomposition, a transit parity, and a characterisation of the
level-two multiplier-weighted Kloosterman-type sum at modulus two.
The first and third are unconditional; the second is an exact
inference from two declared geometric identifications, stated as
explicit hypotheses in the proposition. Nothing in this section
derives any hypothesis of the class of Section~\ref{sec:class}, and
the relation of the present results to that class is stated there.

\begin{definition}[The fibre algebra and its two
four-groups]\label{def:fibrealg}
Let $\mathbb{H}$ denote the quaternion algebra over $\R$ with basis
$1, i, j, k$ and relations $i^2 = j^2 = k^2 = ijk = -1$
\cite{ConwaySmith}, and let
\[
Q_8 \;=\; \{\pm 1,\ \pm i,\ \pm j,\ \pm k\} \;\subset\;
\mathbb{H}^{\times}
\]
be the quaternion group, with centre $Z = \{\pm 1\} = \langle -1
\rangle$. Two Klein four-groups act in what follows. The first is
geometric: the two-torsion subgroup $\T^2[2] \subset \T^2$, whose
underlying set is the fixed-point set of Lemma~\ref{lem:corners},
acting on the torus by the half-period translations. The second is
algebraic: the inner automorphism group
$V = \mathrm{Inn}(Q_8) \cong Q_8/Z \cong \Z/2 \times \Z/2$,
generated by conjugation by $i$ and conjugation by $j$, acting on
the set $Q_8$.
\end{definition}

The two four-groups are isomorphic as abstract groups but act on
different objects, the first on the base torus and the second on the
fibre algebra; no identification between them is made anywhere in
this paper, and the two actions enter the argument in logically
separate roles, the translations through the covering geometry of
the torus and the automorphisms through the orbit structure of
$Q_8$.

\begin{lemma}[Orbit decomposition]\label{lem:orbits}
The action of $V \times Z$ on $Q_8$ --- inner automorphisms together
with multiplication by the centre --- has exactly four orbits,
\[
\{\pm 1\}, \qquad \{\pm i\}, \qquad \{\pm j\}, \qquad \{\pm k\},
\]
each of cardinality two. The centre factor is essential: the inner
automorphisms fix $+1$ and $-1$ separately, as does quaternionic
conjugation $q \mapsto \bar q$, so no group of automorphisms or
anti-automorphisms alone connects the orbit $\{\pm 1\}$.
\end{lemma}

\begin{proof}
Conjugation by $i$ fixes $\pm 1$ and $\pm i$ and sends
$j \mapsto -j$ and $k \mapsto -k$; conjugation by $j$ fixes
$\pm 1$ and $\pm j$ and sends $i \mapsto -i$ and $k \mapsto -k$.
Each of the four displayed pairs is therefore closed under both
generating automorphisms, and multiplication by $-1$ exchanges the
two elements of each pair, so each pair is a single orbit of
$V \times Z$; the pairs are disjoint and cover $Q_8$. The smallest
case, worked by hand: the orbit of $+1$. Every inner automorphism
fixes the centre pointwise, and conjugation fixes $+1$ and $-1$
separately, so neither moves $+1$; the sole generator connecting
$+1$ to $-1$ is the centre itself, and the orbit is $\{\pm 1\}$, of
cardinality two. The full decomposition --- closure under the
generators, connectedness within each pair, disjointness, covering,
and the four cardinalities --- is machine-checked by exhaustive
decision on an explicit eight-element model in
Appendix~\ref{app:lean} (Certificate~8).
\end{proof}

\begin{proposition}[Transit parity]\label{prop:parity}
Every pure imaginary unit quaternion squares to $-1$:
\[
(\pm i)^2 \;=\; (\pm j)^2 \;=\; (\pm k)^2 \;=\; -1 .
\]
Consequently, suppose the following two identifications are
declared: \textup{(i)} each half-period translation of $\T^2$ lifts
to the fibre algebra as right-multiplication by a pure imaginary
unit quaternion; \textup{(ii)} one full covering period of $\T^2$
comprises two half-period translations. Then the cumulative fibre
phase over $n$ full covering periods equals $(-1)^n$ for every
$n \ge 0$, independently of the choice of unit in \textup{(i)}.
\end{proposition}

\begin{proof}
The squares are immediate from the defining relations: $i^2 = -1$,
and $(-i)^2 = (-1)^2 i^2 = -1$; likewise for $j$ and $k$. For the
parity, let $u$ be the unit of identification (i). By (ii) one
covering period multiplies the cumulative phase by $u^2 = -1$, so
over $n$ periods the phase is $(-1)^n$; the value does not depend on
which of the six units is chosen, since all six square to $-1$. The
smallest case, worked by hand: $n = 1$ gives one covering period,
hence two half-period lifts, hence phase $u \cdot u = -1 = (-1)^1$.
The implication is machine-checked with the identifications (i)
and (ii) encoded as named hypothesis fields
(Appendix~\ref{app:lean}, Certificate~8), in the same pattern as
Certificate~3: what is proved is exactly the inference from the
declared identifications, whose truth for any given reduction is a
question about that reduction and is outside the scope of this
paper.
\end{proof}

We turn to the arithmetic object that carries the same parity.

\begin{definition}\label{def:ksum}
Let $\nu$ be a multiplier system of weight one-half on
$\Gamma_0(2)$, in the Rademacher--Zuckerman normalisation
\cite{RademacherZuckerman}. For integers $c \ge 1$ and each $h$ with
$0 \le h < c$ and $\gcd(h,c) = 1$, let $\gamma_{h,c}$ denote the
group element of that normalisation with lower-left entry $c$, and
let $\bar h$ satisfy $h\bar h \equiv -1 \pmod c$. The
multiplier-weighted Kloosterman-type sum is
\begin{equation}\label{eq:ksumdef}
K_2(n, m;\, c) \;=\;
\sum_{\substack{0 \le h < c \\ \gcd(h,c)=1}}
\overline{\nu}(\gamma_{h,c})\,
e^{2\pi i (n \bar h + m h)/c} .
\end{equation}
\end{definition}

The sums \eqref{eq:ksumdef} are not the classical Kloosterman sums
$S(n,m;c) = \sum_{x \in (\Z/c\Z)^{\times}} e^{2\pi i (nx + m\bar
x)/c}$ \cite{Kloosterman}, which are real-valued for all arguments,
the terms at $x$ and $-x$ being complex conjugates; the multiplier
weights $\overline{\nu}(\gamma_{h,c})$ produce genuinely complex
values in general, and the distinction is quantitative at the
modulus of interest, as Proposition~\ref{prop:altchar} records.
Sums of the form \eqref{eq:ksumdef} are the arithmetic companions of
the Fourier coefficients of series of Appell type in exact
expansions of Rademacher type \cite{RademacherZuckerman,
BringmannOno}.

\begin{lemma}\label{lem:ksumtwo}
At $c = 2$ the sum \eqref{eq:ksumdef} with $m = -1$ has the single
term $h = 1$, with $\bar h = 1$, and evaluates to
\begin{equation}\label{eq:ksumtwo}
K_2(n, -1;\, 2) \;=\;
\overline{\nu}(\gamma_{1,2})\; e^{\pi i (n-1)} \;=\;
\overline{\nu}(\gamma_{1,2})\,(-1)^{n-1} .
\end{equation}
\end{lemma}

\begin{proof}
The units modulo $2$ are $\{1\}$, so the sum has the single term
$h = 1$; the congruence $h\bar h \equiv -1 \pmod 2$ gives
$\bar h \equiv 1$. The exponential is
$e^{2\pi i (n \cdot 1 + (-1) \cdot 1)/2} = e^{\pi i (n-1)} =
(-1)^{n-1}$. The smallest cases, worked by hand: $n = 1$ gives
$K_2(1,-1;2) = \overline{\nu}(\gamma_{1,2})$, and $n = 2$ gives
$K_2(2,-1;2) = -\overline{\nu}(\gamma_{1,2})$.
\end{proof}

\begin{proposition}[Characterisation of the
alternation]\label{prop:altchar}
The identity
\begin{equation}\label{eq:altchar}
K_2(n, -1;\, 2) \;=\; (-1)^n \qquad \text{for all } n \ge 1
\end{equation}
holds if and only if $\overline{\nu}(\gamma_{1,2}) = -1$. For the
trivial multiplier --- the classical sum --- the value is
$(-1)^{n-1} = -(-1)^n$, and \eqref{eq:altchar} fails at every $n$.
\end{proposition}

\begin{proof}
By Lemma~\ref{lem:ksumtwo},
$K_2(n,-1;2) = \overline{\nu}(\gamma_{1,2})(-1)^{n-1}$. If
$\overline{\nu}(\gamma_{1,2}) = -1$ then
$K_2(n,-1;2) = -(-1)^{n-1} = (-1)^n$ for every $n$; conversely,
evaluating \eqref{eq:altchar} at $n = 1$ gives
$\overline{\nu}(\gamma_{1,2}) = -1$. Setting
$\overline{\nu} \equiv 1$ in \eqref{eq:ksumtwo} gives the classical
evaluation. Both the characterisation and the classical contrast
are machine-checked (Appendix~\ref{app:lean}, Certificate~8).
\end{proof}

\begin{remark}\label{rem:mult}
We do not evaluate $\overline{\nu}(\gamma_{1,2})$ here.
Proposition~\ref{prop:altchar} states exactly what the parity of
Proposition~\ref{prop:parity} determines and what it does not: the
transit parity $(-1)^n$ coincides with the value of the
multiplier-weighted sum \eqref{eq:ksumdef} at modulus two for all
$n$ precisely when the single multiplier value is $-1$, and we
identify the two under that condition, as an identification rather
than a theorem. The contrast with the classical sum shows that the
identification genuinely selects the multiplier-weighted object:
the two differ by an overall sign at every $n$.
\end{remark}

The relation of this section to the class of
Section~\ref{sec:class} is one of structure, not derivation, and is
stated with the class hypotheses. Here we record only the
observation that the parity of Proposition~\ref{prop:parity} and
the counting function $\dodd$ of Section~\ref{sec:arithmetic} are
attached to the same level-two object: the sums \eqref{eq:ksumdef}
accompany the coefficients of Appell-type series, of which the
level-two sum \eqref{eq:appell} is the member at hand, in
expansions of Rademacher type \cite{BringmannOno}, so that the two
are faces of a single arithmetic structure rather than independent
inputs.

